\documentclass[12pt]{article}
\usepackage[utf8]{inputenc}
\usepackage[T1]{fontenc}
\usepackage{hyperref}
\title{Watakataui Language Resources for K-12 Curriculum Development}
\author{Compiled by Claude Sonnet 5.5 (claude-sonnet-5-5) for Andr\'as Kornai}
\date{}
\begin{document}
\maketitle

\section{Introduction}
Watakataui (ISO 639-3 wtk; Glottocode wata1251) is an endangered Sepik language of East Sepik Province, Papua New Guinea. Search results place it in the Bahinemo grouping and in two villages. Speaker counts range from about 160 to 930, depending on the source, and the figures do not come from one consistent survey. Alternate names include Waxe and Waxei.

\section{Primary Educational Resources}
I found no primer, textbook or school programme for Watakataui.

\section{Linguistic and Reference Materials}
\begin{itemize}
\item Overview: \url{https://en.wikipedia.org/wiki/Watakataui_language}
\item Verbix endangered languages page: \url{https://docs.verbix.com/EndangeredLanguages/Watakataui}
\item ASJP word list entry: \url{https://asjp.clld.org/languages/WATAKAKAUI}
\item Crossgram entry: \url{https://crossgram.clld.org/languages/wata1251}
\item Endangered Languages Project entry: \url{https://www.endangeredlanguages.com/lang/10618?hl=en}
\item Glottolog bibliography entries for Yoichi Yamada (1995, 1998): \url{https://glottolog.org/resource/reference/id/13246} and \url{https://glottolog.org/resource/reference/id/73286}
\end{itemize}
I did not check how much linguistic detail these entries hold.

\section{Religious and Community Resources}
Joshua Project reports Bible portions (1991) and a New Testament (2014) in the language: \url{https://joshuaproject.net/index.php/languages/wtk}. I did not check this against a publisher's record.

\section{Digital and Audio Resources}
Global Recordings Network has a page for this language code (\url{https://globalrecordings.net/en/language/wtk}). The search summary reported no recordings listed there.

\section{Curriculum Development Context}
The language is not known to be taught in schools, and I found no policy documents about it. Any curriculum would have to be built with the community, starting from documentation and the scripture translation.

\end{document}